\chapter{Développer du code à visée de production dans un existant endetté}
\chaptermark{Développer du code à visée de production}
\label{chap:production}
% Chapitre-chapeau du volet 3 : fonctionnalité d'immuabilité (t7).

Ce dernier chapitre, consacré à une mission menée chez Beemo, change de
posture : il ne s'agit plus
d'outiller ou de tester l'existant, mais d'\textbf{écrire du code neuf
destiné à la production} au cœur d'un firmware massivement endetté : fichiers
de plusieurs dizaines de milliers de lignes, couplage généralisé, absence de
tests (chapitre~\ref{chap:veille}).

L'enjeu est de livrer la fonctionnalité \textbf{sans accroître la dette}, et
la difficulté est double. Technique, d'abord : chaque point de contact avec
le legacy importe son opacité et son couplage ; du code neuf écrit « dans le
style » de l'existant en hérite mécaniquement, et l'absence de filet de test
rend toute intégration risquée. Méthodologique, ensuite : l'existant exerce
une pression de conformité, car le chemin de moindre effort consiste à imiter
ce qui est déjà là, donc à reproduire les pratiques qui ont créé la dette.
S'en prémunir suppose une discipline d'\textbf{isolation} : écrire le neuf
de manière testable et découplée, et confiner le dialogue avec le legacy dans
une couche d'adaptation explicite, plutôt que de laisser les deux mondes
s'entremêler.

La section~\ref{chap:t7} illustre cette démarche avec le développement d'une
fonctionnalité d'immuabilité des sauvegardes, dont l'architecture trace une
voie de modernisation progressive du firmware.
